\section{Finite Fourier selection and exactness}
\label{sec:fourier-selection}

The construction begins with a finite-group coefficient. We record the selection identity in a form suited to a fixed rational target and keep the later passage from a quotient congruence to exact equality logically separate.

\subsection{Character orthogonality}

Let $A$ be a finite abelian group and let $\widehat A$ be its character group. For $a,t\in A$,
\[
 \frac1{|A|}\sum_{\chi\in\widehat A}\overline{\chi(t)}\chi(a)=\mathbf 1_{a=t}.
\]
Consequently, if $T$ is an $A$-valued random variable, then
\begin{equation}
\label{eq:finite-fourier-coefficient}
 \Prob(T=t)=\frac1{|A|}\sum_{\chi\in\widehat A}\overline{\chi(t)}\E\chi(T).
\end{equation}
The numerator is the real nonnegative number $|A|\Prob(T=t)$, so it is enough to prove that its real part is positive.

\begin{lemma}[Independent product formula]
\label{lem:independent-product}
Let $v_j\in A$, and let $\xi_j$ be independent Bernoulli random variables with parameters $\theta_j\in[0,1]$. For $T=\sum_j\xi_jv_j$ and every $\chi\in\widehat A$,
\[
 \E\chi(T)=\prod_j\bigl((1-\theta_j)+\theta_j\chi(v_j)\bigr).
\]
\end{lemma}
\begin{proof}
Since $\chi$ is a homomorphism, $\chi(T)=\prod_j\chi(v_j)^{\xi_j}$. Independence gives the product of expectations.
\end{proof}

\subsection{Reciprocal subset sums at a fixed target}

Let $\mathcal E$ be a finite set of positive integers and let $L$ be divisible by every $e\in\mathcal E$. Associate to $e$ the residue $v_e=L/e\pmod L$. Let
\[
 t=\frac ab\in(0,1)
\]
be reduced, with $b\mid L$, and let $\xi_e$ be independent Bernoulli variables with common parameter $\theta$. Then
\[
 T=\sum_{e\in\mathcal E}\xi_ev_e
\]
encodes the subset reciprocal sum modulo one: the congruence $T=aL/b\pmod L$ is equivalent to
\[
 \sum_{e\in\mathcal E}\frac{\xi_e}{e}\equiv\frac ab\pmod1.
\]
Writing characters as $h\mapsto\exp(2\pi i h\,\cdot/L)$, \cref{eq:finite-fourier-coefficient} gives
\begin{equation}
\label{eq:reciprocal-fourier}
 L\Prob\left(T=\frac{aL}{b}\right)=\sum_{h\bmod L}\exp\left(-\frac{2\pi i a h}{b}\right)\prod_{e\in\mathcal E}\left((1-\theta)+\theta\exp\left(\frac{2\pi i h}{e}\right)\right).
\end{equation}
We shall prove that the real part of the right side is positive.

The linear phase is centred by choosing
\begin{equation}
\label{eq:theta-centering}
 \theta\sum_{e\in\mathcal E}\frac1e=t.
\end{equation}
For the parameterized semiprime family this is $\theta=t/\Lambda$, where $\Lambda$ is the total reciprocal load. Apart from its scalar effect through this centring parameter, the numerator $a$ enters the Fourier geometry only through the target character in \cref{eq:reciprocal-fourier}; the CRT rows which observe the target depend only on the prime divisors of $b$.

\subsection{From congruence to equality}

The Fourier coefficient gives a reciprocal subset sum congruent to $t$ modulo one. In the present construction every subset sum lies in $[0,\Lambda]$ and $\Lambda<1$, while $t\in(0,1)$. Hence the congruence can represent only exact equality. We shall invoke this no-wrap observation only after positivity has been proved.

\subsection{Chinese-remainder coordinates}

The period $L$ is squarefree. Hence the Chinese remainder theorem identifies
\[
 \mathbb Z/L\mathbb Z\cong\prod_{p\mid L}\mathbb Z/p\mathbb Z.
\]
A chosen block of large prime coordinates forms the anchor block. The remaining prime coordinates are grouped into anchor-cross rows whose kernels depend on the anchor assignment. Denominator factors involving two anchor primes contribute to the anchor weight; factors involving one anchor prime and one row prime contribute to that row; every other factor remains in a residual term after the row choices have been compressed. This is an exact partition of the original Fourier product.
